% =====================================================================
\section{Seminarios prácticos}
% =====================================================================
Esta sección resuelve los dieciséis problemas de la lista del tema \cite{gonzalez1} y la cuestión 1 del
parcial de 2025/26 \cite{parcial}. Cada problema sigue la misma estructura: el \emph{enunciado}, la
\emph{idea} que guía la resolución, la \emph{resolución} paso a paso con todo el desarrollo simbólico
y, cuando aporta, una \emph{observación} sobre el resultado y una figura. Salvo indicación contraria,
$R=\SI{8,314}{J/(mol.K)}$.

\subsection{Cómo plantear un problema del Tema 1}

Todos los problemas del tema se reducen a cinco tareas (tabla~\ref{tab:metodos}). La primera decisión
es identificar \emph{qué se da} (ecuación de estado, coeficientes, condiciones de equilibrio o una
relación de proceso) y \emph{qué se pide}.

\begin{table}[htbp]
\centering
\caption{Mapa de decisión: qué método conviene según lo que pide el enunciado.}
\label{tab:metodos}
\small
\begin{tabularx}{\textwidth}{Xlr}
\toprule
Lo que da / pide el enunciado & Método & Problemas\\
\midrule
Condiciones de equilibrio térmico entre sistemas; hallar $\theta_i$ & M1. Ley cero & 1.1\\
Proceso $pV^n=\cte$; hallar $p_2$ o el exponente & M2. Politrópico & 1.3, 1.4\\
Dilatación, calentamiento a $V$ cte, cambios pequeños & M3. Coeficientes constantes & 1.5--1.7, 1.16\\
Ecuación de estado dada; hallar $\alpha$, $\kappa_T$, $\beta$ & M4. Derivación & 1.8, P.1(b)\\
Coeficientes (o $\dd V$) dados; hallar la ecuación de estado & M5. Integración & 1.2, 1.9--1.15, P.1(a)\\
\bottomrule
\end{tabularx}
\end{table}

\begin{recuadro}{Los cinco métodos}
\begin{enumerate}[label=\textbf{\color{navy}M\arabic*.},leftmargin=2.2em,itemsep=4pt]
\item \emph{Ley cero.} Reescribir cada condición de equilibrio aislando la misma combinación
  $\theta_3(p_3,V_3)$ del sistema común; lo que queda a cada lado son $\theta_1$ y $\theta_2$, y
  $\Sigma_1\sim\Sigma_2\iff\theta_1=\theta_2$.
\item \emph{Politrópico.} $p_1V_1^n=p_2V_2^n\Rightarrow p_2=p_1(V_1/V_2)^n$ y
  $n=\ln(p_2/p_1)/\ln(V_1/V_2)$. Unificar unidades \emph{antes} de operar.
\item \emph{Coeficientes constantes.} $\Delta L\simeq\alpha_lL_0\Delta T$;
  $\ln(V_2/V_1)=\alpha\Delta T-\kappa_T\Delta p$; a volumen fijo, $\Delta p=(\alpha/\kappa_T)\Delta T$.
\item \emph{Derivación.} Despejar la variable que se deriva si es sencillo; si no, usar la regla
  recíproca \eqref{eq:reciclica} o derivar implícitamente. Expresar el resultado en las variables pedidas.
\item \emph{Integración.} Escribir la diferencial adecuada \eqref{eq:resumen}, comprobar su exactitud
  \eqref{eq:exacta}, integrar y fijar la constante con un límite físico. Es el método más largo; su
  receta completa, en ocho pasos, se da a continuación.
\end{enumerate}
\end{recuadro}

\subsection{El método M5 paso a paso}\label{sec:receta}

Los problemas~1.2 y 1.9--1.15 y la cuestión P.1(a) se resuelven \emph{siempre} con la misma secuencia.
Escribirla explícitamente en el examen es lo que demuestra que se sabe resolver el problema, y no
solo que se ha llegado al resultado.

\paso{Paso 0: planteamiento} Identificar la \emph{incógnita} como función de las otras dos variables
---normalmente $\Vm=\Vm(T,p)$, o bien $p=p(T,\Vm)$ si así lo pide el enunciado---, los \emph{datos}
(dos de los tres coeficientes) y las \emph{condiciones} (un límite físico o un estado de referencia).

\paso{Paso 0 bis: completar el par} Si los datos no son los coeficientes que aparecen en la
diferencial de la incógnita, el que falta se obtiene de la relación cíclica \eqref{eq:apbk}:
\[
  \alpha=p\,\beta\,\kappa_T,\qquad \kappa_T=\frac{\alpha}{p\beta},\qquad \beta=\frac{\alpha}{p\kappa_T}.
\]

\paso{Paso 1: diferencial total} Escribir la diferencial de la incógnita y sustituir cada derivada por
su coeficiente con \eqref{eq:derivadas}--\eqref{eq:inversas}:
\[
  \dd\Vm=\Vm\alpha\,\dd T-\Vm\kappa_T\,\dd p\qquad\text{o}\qquad
  \dd p=p\beta\,\dd T-\frac{\dd\Vm}{\Vm\kappa_T}.
\]

\paso{Paso 2: elegir la forma} Se elige la forma en la que \emph{la incógnita desaparece de los
coeficientes}, porque solo entonces la diferencial se puede integrar en las otras dos variables
(tabla~\ref{tab:forma}).

\begin{table}[htbp]
\centering
\caption{Criterio para elegir la forma de la diferencial en el paso 2.}
\label{tab:forma}
\small
\renewcommand{\arraystretch}{1.6}
\begin{tabular}{lll}
\toprule
Los coeficientes $\alpha$, $\kappa_T$\dots & Forma que conviene & Problemas\\
\midrule
no contienen la incógnita $X$ & $\dd\ln X=\dfrac{\dd X}{X}$ & 1.9, 1.12, P.1\\[4pt]
contienen el factor $1/X$ & $\dd X$ (multiplicar por $X$) & 1.2, 1.10, 1.11\\[4pt]
contienen el factor $(X-A)/X$ & $\dd\ln(X-A)=\dfrac{\dd X}{X-A}$ & 1.13\\[4pt]
$p\kappa_T$ no contiene $p$ & $\dd\ln p=\beta\,\dd T-\dfrac{\dd\Vm}{p\kappa_T\Vm}$ & 1.14\\
\bottomrule
\end{tabular}
\end{table}

\paso{Paso 3: sustituir y simplificar} Sustituir los datos y escribir el resultado como
$\dd z=M(x,y)\,\dd x+N(x,y)\,\dd y$, separando cada coeficiente en fracciones simples
(por ejemplo, $\frac{RT\Vm+a}{RT^2\Vm}=\frac1T+\frac{a}{RT^2\Vm}$).

\paso{Paso 4: exactitud} Comprobar \eqref{eq:exacta}, $\partial_yM=\partial_xN$. Para derivar las
fracciones se usa la regla del cociente,
\begin{equation}\label{eq:cociente}
  \left(\frac fg\right)'=\frac{f'g-fg'}{g^2}
  \qquad\xrightarrow{\;f=\text{cte}\;}\qquad
  \left(\frac fg\right)'=-\frac{f\,g'}{g^2},
\end{equation}
que, con numerador constante, coincide con la regla de la cadena aplicada a $f\,g^{-1}$. Conviene
escribir $g'$ entre paréntesis con su signo antes de multiplicar.

\paso{Paso 5: integrar} Se distinguen dos casos:
\begin{itemize}[leftmargin=1.6em]
\item[\textbf{\color{navy}A.}] \emph{Variables separadas}: $M=M(x)$ y $N=N(y)$. Se integra cada término
  por su lado: $z=\int M\,\dd x+\int N\,\dd y+\cte$.
\item[\textbf{\color{navy}B.}] \emph{Variables no separadas}. (5.1) Se integra $M$ en $x$ con $y$ fija;
  la «constante» depende de $y$: $z=\int M\,\dd x+g(y)$. (5.2) Se deriva respecto de $y$ y se iguala a
  $N$, lo que da $g'(y)$; los términos en $x$ \emph{deben} cancelarse (si no, hay un error, porque la
  forma es exacta). (5.3) Se integra $g'(y)$.
\end{itemize}
Las primitivas que aparecen en el tema son siempre de tres tipos:
\begin{equation}\label{eq:primitivas}
  \int\frac{f'(x)}{f(x)}\,\dd x=\ln f(x),\qquad
  \int\frac{a\,\dd x}{c+ax}=\ln(c+ax),\qquad
  \int x^{-n}\,\dd x=\frac{x^{1-n}}{1-n}\ (n\neq1),
\end{equation}
donde en la segunda $c$ puede contener la otra variable, que está fija (cambio $u=c+ax$, $\dd u=a\,\dd x$).

\paso{Paso 6: quitar logaritmos y fijar la constante} Conviene escribir la constante de integración
directamente como $\ln K$ y agrupar con
\begin{equation}\label{eq:logs}
  a\ln x=\ln x^a,\quad\ \ln x\pm\ln y=\ln(xy^{\pm1}),\quad\ 
  \ln X=\ln Y\iff X=Y,\quad\ \e^{\ln X+y}=X\,\e^{y}.
\end{equation}
Si se escribió $+C$, basta tomar exponenciales: $\e^{\ln X+C}=\e^{C}X$, y $K\equiv\e^C$ es otra
constante. La constante se fija con el límite físico del enunciado (normalmente gas ideal para
$\Vm\to\infty$); si no hay ninguno, queda libre.

\paso{Paso 7: comprobación} Derivar el resultado y recuperar los coeficientes del enunciado. Es la
prueba definitiva de que no hay errores de signo.

\begin{figure}[htbp]
\centering
\begin{tikzpicture}[font=\footnotesize,
   p/.style={draw=navy,fill=fondo,align=center,rounded corners=1pt,inner sep=3pt,text width=2.6cm,minimum height=1.0cm},
   d/.style={draw=acento,fill=white,diamond,aspect=2.4,align=center,inner sep=1pt,font=\scriptsize},
   f/.style={-{Stealth[length=4pt]},navy}]
  \foreach \n/\x/\t in {s0/0/{\textbf{0} Planteamiento\\ incógnita, datos},
                        s1/3.15/{\textbf{0 bis, 1} Par de\\ coef.; diferencial},
                        s2/6.3/{\textbf{2} Forma\\ $\dd X$ o $\dd\ln X$},
                        s3/9.45/{\textbf{3} Sustituir\\ $M\,\dd x+N\,\dd y$},
                        s4/12.6/{\textbf{4} Exactitud\\ $\partial_yM=\partial_xN$}}
     \node[p] (\n) at (\x,0) {\t};
  \node[d] (dd) at (12.6,-2.25) {¿separadas?};
  \node[p] (sA) at (9.45,-1.6) {\textbf{5A} Integrar\\ cada término};
  \node[p] (sB) at (9.45,-2.9) {\textbf{5B} $\int\!M\dd x+g(y)$,\\ $g$ con $N$};
  \node[p] (s6) at (6.3,-2.25) {\textbf{6} Quitar $\ln$;\\ fijar constante};
  \node[p] (s7) at (3.15,-2.25) {\textbf{7} Comprobar\\ $\alpha$, $\kappa_T$, $\beta$};
  \draw[f] (s0) -- (s1); \draw[f] (s1) -- (s2); \draw[f] (s2) -- (s3); \draw[f] (s3) -- (s4);
  \draw[f] (s4) -- (dd);
  \draw[f] (dd.west) -- node[above,font=\scriptsize,acento] {sí} (sA.east);
  \draw[f] (dd.west) -- node[below,font=\scriptsize,acento] {no} (sB.east);
  \draw[f] (sA.west) -- (s6.east);
  \draw[f] (sB.west) -- (s6.east);
  \draw[f] (s6) -- (s7);
\end{tikzpicture}
\caption{Esquema del método M5. La única bifurcación es si las variables están separadas (caso A) o no
(caso B, con la función auxiliar $g$ de la otra variable).}
\label{fig:receta}
\end{figure}

\begin{observacion}[Atajo de la diferencial logarítmica]
Si los coeficientes de $\dd T$ y de $\dd p$ tienen el \emph{mismo denominador} $D(T,p)$ y sus numeradores
son $\partial_TD$ y $\partial_pD$, la forma completa es $\dd D/D=\dd\ln D$ y se integra en una línea
(problemas~1.12, 1.13 y P.1). Si no se ve el atajo, el caso B lleva al mismo resultado: la función
auxiliar sale constante, $g'=0$ (problema~1.12).
\end{observacion}

% ---------------------------------------------------------------------
\subsection{Problema 1.1. Temperaturas empíricas de tres gases}
% ---------------------------------------------------------------------
\begin{problema}{1.1}{Ley cero con tres gases}
Los sistemas $\Sigma_1$, $\Sigma_2$ y $\Sigma_3$ son gases de variables $(p_i,V_i)$. Cuando $\Sigma_1$ y
$\Sigma_3$ están en equilibrio térmico se cumple $p_1V_1-bp_1-p_3V_3=0$, y cuando lo están $\Sigma_2$ y
$\Sigma_3$, $p_2V_2-p_3V_3+b'p_3V_3/V_2=0$, con $b,b'$ constantes. ¿Qué relación expresa el equilibrio
térmico entre $\Sigma_1$ y $\Sigma_2$? ¿Cuáles son las funciones $\theta_i(p_i,V_i)$ que coinciden en el
equilibrio?
\end{problema}

\begin{resolucion}
\idea Método M1: aislar en ambas condiciones la \emph{misma} combinación de $\Sigma_3$, que aparece de
forma natural como $p_3V_3$.

\paso{Paso 0: planteamiento} Se buscan funciones $\theta_i(p_i,V_i)$ tales que
$\Sigma_i\sim\Sigma_j\iff\theta_i=\theta_j$. Como $\Sigma_3$ es el sistema común, basta despejar en cada
condición la parte que depende solo de $(p_3,V_3)$.

\paso{Paso 1: condición \texorpdfstring{$\Sigma_1\sim\Sigma_3$}{S1~S3}} Sacando factor común $p_1$ y
pasando $p_3V_3$ al otro miembro:
\[
  p_1V_1-bp_1-p_3V_3=0\quad\Longrightarrow\quad p_1(V_1-b)=p_3V_3.
\]

\paso{Paso 2: condición \texorpdfstring{$\Sigma_2\sim\Sigma_3$}{S2~S3}} Se agrupan los dos términos en
$p_3V_3$ y se saca factor común:
\[
  p_2V_2=p_3V_3-\frac{b'p_3V_3}{V_2}=p_3V_3\Bigl(1-\frac{b'}{V_2}\Bigr)=p_3V_3\,\frac{V_2-b'}{V_2}
  \quad\Longrightarrow\quad p_3V_3=\frac{p_2V_2^2}{V_2-b'}.
\]

\paso{Paso 3: ley cero} Si $\Sigma_1\sim\Sigma_3$ y $\Sigma_2\sim\Sigma_3$, por la ley cero
$\Sigma_1\sim\Sigma_2$, y las dos expresiones de $p_3V_3$ deben coincidir:
\[
  \boxed{\,p_1(V_1-b)=\frac{p_2V_2^2}{V_2-b'}\,}
\]
Es la relación pedida: no contiene las variables de $\Sigma_3$, como exige la ley cero.

\paso{Paso 4: temperaturas empíricas} Cada miembro depende solo de las variables de un sistema y los
tres valores coinciden en el equilibrio mutuo, $\theta_1=\theta_3=\theta_2$:
\[
  \boxed{\,\theta_1=p_1(V_1-b),\qquad\theta_2=\frac{p_2V_2^2}{V_2-b'},\qquad\theta_3=p_3V_3\,}
\]

\paso{Paso 5: comprobación} Sustituyendo $\theta_1=\theta_3$ se recupera $p_1V_1-bp_1-p_3V_3=0$; y
sustituyendo $\theta_2=\theta_3$ y multiplicando por $(V_2-b')/V_2$, se recupera
$p_2V_2-p_3V_3+b'p_3V_3/V_2=0$.
\end{resolucion}

\begin{observacion}[Unicidad]
Cualquier función monótona $\phi(\theta)$ aplicada a las tres $\theta_i$ es igualmente válida; la elección
anterior es la más simple. $\Sigma_3$ es un gas ideal ($\theta_3\propto T$), $\Sigma_1$ un gas de Clausius y
$\Sigma_2$ un gas con corrección en $b'/V_2$. La figura~\ref{fig:p11} muestra que la «misma temperatura»
corresponde a isotermas distintas en cada sistema.
\end{observacion}

\begin{figure}[htbp]
\centering
\begin{tikzpicture}
\begin{axis}[graf,leyenda fuera,width=8cm,height=5.6cm,xmin=0.6,xmax=4,ymin=0,ymax=4.2,
    xlabel={$V_i$ (u.\,arb.)},ylabel={$p_i$ (u.\,arb.)}]
  \addplot[navy,domain=0.62:4,samples=80] {1.5/(x-0.3)};
  \addplot[acento,domain=0.62:4,samples=80] {1.5*(x-0.2)/x^2};
  \addplot[rojo,densely dashed,domain=0.62:4,samples=80] {1.5/x};
  \addplot[navy,domain=0.85:4,samples=80,forget plot,opacity=0.45] {3/(x-0.3)};
  \addplot[acento,domain=0.62:4,samples=80,forget plot,opacity=0.45] {3*(x-0.2)/x^2};
  \addplot[rojo,densely dashed,domain=0.72:4,samples=80,forget plot,opacity=0.45] {3/x};
  \legend{$\Sigma_1$: $p_1(V_1-b)=\theta$,$\Sigma_2$: $p_2V_2^2/(V_2-b')=\theta$,$\Sigma_3$: $p_3V_3=\theta$}
\end{axis}
\end{tikzpicture}
\caption{Problema 1.1. Isotermas $\theta_i(p_i,V_i)=\theta$ de los tres gases para $\theta=1{,}5$ (trazo
intenso) y $\theta=3$ (trazo suave), con $b=0{,}3$ y $b'=0{,}2$. Estados de sistemas distintos situados
sobre curvas con el mismo $\theta$ están en equilibrio térmico mutuo.}
\label{fig:p11}
\end{figure}

% ---------------------------------------------------------------------
\subsection{Problema 1.2. Ecuación de estado a partir de \texorpdfstring{$\dd V$}{dV}}
% ---------------------------------------------------------------------
\begin{problema}{1.2}{Integración directa}
Un fluido pasa de $(T,p)$ a $(T+\dd T,p+\dd p)$ con $\displaystyle\dd V=\frac{2V}{T}\,\dd T-\frac{3V}{p}\,\dd p$.
Encuéntrese la relación entre $V$, $p$ y $T$.
\end{problema}

\begin{resolucion}
\idea Método M5. Al comparar con $\dd V=\alpha V\dd T-\kappa_TV\dd p$ se leen los coeficientes; como estos
no contienen $V$, se trabaja con $\dd\ln V$ y las variables quedan separadas.

\paso{Paso 0: planteamiento} Incógnita $V=V(T,p)$. Identificando término a término con
\eqref{eq:dlnV} multiplicada por $V$:
\[
  \alpha V=\frac{2V}{T}\ \Longrightarrow\ \alpha=\frac2T,\qquad
  \kappa_TV=\frac{3V}{p}\ \Longrightarrow\ \kappa_T=\frac3p.
\]

\paso{Paso 1--2: diferencial y forma} Los coeficientes dados llevan el factor $V$; dividiendo entre $V$,
\[
  \frac{\dd V}{V}=\dd\ln V=\frac{2}{T}\,\dd T-\frac{3}{p}\,\dd p\qquad
  \Bigl(M=\frac2T,\quad N=-\frac3p\Bigr).
\]

\paso{Paso 3: exactitud}
\[
  \pd{M}{p}{T}=\frac{\partial}{\partial p}\Bigl(\frac2T\Bigr)=0,\qquad
  \pd{N}{T}{p}=\frac{\partial}{\partial T}\Bigl(-\frac3p\Bigr)=0\quad\Longrightarrow\quad\text{exacta.}
\]

\paso{Paso 4: integración (caso A)} $M$ solo depende de $T$ y $N$ solo de $p$; se integra cada término
por su lado, con la constante escrita como $\ln K$:
\[
  \ln V=2\int\frac{\dd T}{T}-3\int\frac{\dd p}{p}=2\ln T-3\ln p+\ln K.
\]

\paso{Paso 5: quitar logaritmos} Con \eqref{eq:logs}, $2\ln T=\ln T^2$, $3\ln p=\ln p^3$, y
\[
  \ln V=\ln\frac{K\,T^2}{p^3}\quad\Longrightarrow\quad V=\frac{K\,T^2}{p^3}\quad\Longrightarrow\quad
  \boxed{\,p^3V=K\,T^2\,}
\]

\paso{Paso 6: comprobación} Con $V=KT^2p^{-3}$:
\[
  \alpha=\frac1V\cdot2KTp^{-3}=\frac2T\chk\qquad
  \kappa_T=-\frac1V\cdot\bigl(-3KT^2p^{-4}\bigr)=\frac3p\chk
\]
\end{resolucion}

\begin{figure}[htbp]
\centering
\begin{tikzpicture}
\begin{axis}[graf,leyenda fuera,width=8cm,height=5.4cm,xmin=0.2,xmax=4,ymin=0,ymax=3.2,
    xlabel={$V$ (u.\,arb.)},ylabel={$p$ (u.\,arb.)}]
  \addplot[acentoclaro,domain=0.2:4,samples=80] {(1/x)^(1/3)};
  \addplot[acentomedio,domain=0.2:4,samples=80] {(2.25/x)^(1/3)};
  \addplot[navy,domain=0.2:4,samples=80] {(4/x)^(1/3)};
  \addplot[rojo,densely dashed,domain=0.2:4,samples=80] {(1/x)};
  \legend{$T=1$,$T=1{,}5$,$T=2$,gas ideal ($T=1$)}
\end{axis}
\end{tikzpicture}
\caption{Problema 1.2. Isotermas $p=(KT^2/V)^{1/3}$ ($K=1$): decaen como $V^{-1/3}$, mucho más
despacio que las del gas ideal ($V^{-1}$), porque $\kappa_T=3/p$ es el triple que el del gas ideal.}
\label{fig:p12}
\end{figure}

% ---------------------------------------------------------------------
\subsection{Problema 1.3. Proceso politrópico: presión final}
% ---------------------------------------------------------------------
\begin{problema}{1.3}{Expansión con $pV^{7/5}=\cte$}
Un gas efectúa un proceso reversible con $pV^{7/5}=\cte$ desde $p_1=\SI{1}{bar}$, $V_1=\SI{1}{m^3}$
hasta $V_2=\SI{3e6}{cm^3}$. Calcúlese $p_2$ en kPa y represéntese el proceso en el diagrama $pV$.
\end{problema}

\begin{resolucion}
\idea Método M2 tras unificar unidades.

\paso{Paso 0: unidades} Con $\SI{1}{m^3}=\SI{e6}{cm^3}$ y $\SI{1}{bar}=\SI{e5}{Pa}=\SI{100}{kPa}$:
\[
  V_2=\SI{3e6}{cm^3}\cdot\frac{\SI{1}{m^3}}{\SI{e6}{cm^3}}=\SI{3}{m^3},\qquad p_1=\SI{100}{kPa},\qquad
  \frac{V_1}{V_2}=\frac13.
\]

\paso{Paso 1: ecuación del proceso} Como $pV^n$ toma el mismo valor en todos los estados del proceso,
con $n=7/5=1{,}4$:
\[
  p_1V_1^{\,n}=p_2V_2^{\,n}\quad\Longrightarrow\quad p_2=p_1\Bigl(\frac{V_1}{V_2}\Bigr)^{n}.
\]

\paso{Paso 2: cálculo} La potencia se evalúa con logaritmos:
\[
  \Bigl(\frac13\Bigr)^{1{,}4}=\e^{-1{,}4\ln3}=\e^{-1{,}5381}=\frac{1}{4{,}656}=0{,}2148
  \quad\Longrightarrow\quad p_2=\SI{100}{kPa}\cdot0{,}2148\simeq\boxed{\,\SI{21,5}{kPa}\,}
\]

\paso{Paso 3: representación} Por el estado inicial pasan la politrópica $p=p_1(V_1/V)^{1{,}4}$ y la
isoterma $p=p_1V_1/V$ (que llegaría a $p_1/3=\SI{33,3}{kPa}$). Como $n>1$ y $V_1/V<1$,
$(V_1/V)^{1{,}4}<V_1/V$: la politrópica cae \emph{más deprisa} que la isoterma (figura~\ref{fig:p13}).
\end{resolucion}

\begin{figure}[htbp]
\centering
\begin{tikzpicture}
\begin{axis}[graf,leyenda fuera,width=7.6cm,height=5.4cm,xmin=0.8,xmax=3.2,ymin=0,ymax=110,
    xtick={1,2,3},ytick={21.5,33.3,100},yticklabels={21{,}5,33{,}3,100},
    xlabel={$V$ (m$^3$)},ylabel={$p$ (kPa)}]
  \addplot[draw=none,fill=fondomedio,domain=1:3,samples=60,forget plot] {100*x^(-1.4)} \closedcycle;
  \addplot[navy,domain=1:3,samples=60] {100*x^(-1.4)};
  \addplot[acento,densely dashed,domain=1:3,samples=60] {100/x};
  \fill[rojo] (axis cs:1,100) circle (1.8pt) node[above right,font=\scriptsize,black] {1};
  \fill[rojo] (axis cs:3,21.48) circle (1.8pt) node[below left,font=\scriptsize,black] {2};
  \draw[navy,-{Stealth[length=5pt]}] (axis cs:1.35,80) -- (axis cs:1.75,60);
  \legend{$pV^{7/5}=\cte$,isoterma $pV=\cte$}
\end{axis}
\end{tikzpicture}
\caption{Problema 1.3. Expansión politrópica con $n=7/5$ (continua) frente a la isoterma por el estado
inicial (discontinua). El área sombreada es el trabajo realizado por el gas, que se estudiará en el Tema 2.}
\label{fig:p13}
\end{figure}

% ---------------------------------------------------------------------
\subsection{Problema 1.4. Exponente politrópico}
% ---------------------------------------------------------------------
\begin{problema}{1.4}{Exponente de una compresión}
Un cilindro--pistón contiene un gas que se comprime reversiblemente desde $p_1=\SI{2}{bar}$,
$v_1=\SI{83,54}{cm^3/kg}$ hasta $p_2=\SI{10}{bar}$, $v_2=\SI{21,34}{cm^3/kg}$, según $pv^\alpha=\cte$.
Calcúlese $\alpha$.
\end{problema}

\begin{resolucion}
\idea Método M2: tomando logaritmos, $pv^\alpha=\cte$ es una recta de pendiente $-\alpha$ en el plano
$(\ln v,\ln p)$. Solo intervienen cocientes, así que no hace falta convertir unidades.

\paso{Paso 1: ecuación del proceso} Entre los dos estados,
\[
  p_1v_1^{\,\alpha}=p_2v_2^{\,\alpha}\quad\Longrightarrow\quad\frac{p_2}{p_1}=\Bigl(\frac{v_1}{v_2}\Bigr)^{\alpha}.
\]

\paso{Paso 2: despeje} Tomando logaritmos y usando $\ln x^\alpha=\alpha\ln x$:
\[
  \ln\frac{p_2}{p_1}=\alpha\ln\frac{v_1}{v_2}\quad\Longrightarrow\quad
  \alpha=\frac{\ln(p_2/p_1)}{\ln(v_1/v_2)}.
\]

\paso{Paso 3: cálculo} Con $p_2/p_1=5$ y $v_1/v_2=83{,}54/21{,}34=3{,}9147$:
\[
  \alpha=\frac{\ln5}{\ln3{,}9147}=\frac{1{,}6094}{1{,}3648}\quad\Longrightarrow\quad\boxed{\,\alpha\simeq1{,}18\,}
\]

\paso{Paso 4: comprobación} $5^{1/1{,}18}\simeq\e^{1{,}6094/1{,}18}=\e^{1{,}364}\simeq3{,}91=v_1/v_2$\chk\
Como $1<\alpha<\gamma\simeq1{,}4$, la compresión está entre la isoterma y la adiabática del gas ideal diatómico.
\end{resolucion}

\begin{figure}[htbp]
\centering
\begin{tikzpicture}
\begin{axis}[graf,width=7.2cm,height=5.2cm,xmode=log,ymode=log,xmin=18,xmax=100,ymin=1.5,ymax=13,
    xtick={21.34,83.54},xticklabels={21{,}34,83{,}54},ytick={2,10},yticklabels={2,10},
    minor tick num=0,log basis x=10,xlabel={$v$ (cm$^3$/kg)},ylabel={$p$ (bar)}]
  \addplot[navy,domain=21.34:83.54,samples=2] {2*(83.54/x)^1.1793};
  \fill[rojo] (axis cs:83.54,2) circle (1.8pt) node[below,font=\scriptsize,black] {1};
  \fill[rojo] (axis cs:21.34,10) circle (1.8pt) node[above,font=\scriptsize,black] {2};
  \draw[navy,-{Stealth[length=5pt]}] (axis cs:55,6) -- (axis cs:38,7.6) node[pos=0,right,font=\scriptsize,acento] {compresión};
  \node[acento,font=\scriptsize] at (axis cs:32,2.0) {pendiente $=-\alpha\simeq-1{,}18$};
\end{axis}
\end{tikzpicture}
\caption{Problema 1.4. En ejes doblemente logarítmicos un proceso politrópico es una recta; su
pendiente es el exponente cambiado de signo.}
\label{fig:p14}
\end{figure}

% ---------------------------------------------------------------------
\subsection{Problema 1.5. Reloj de péndulo}
% ---------------------------------------------------------------------
\begin{problema}{1.5}{Péndulo de cobre}
Un reloj con un péndulo de cobre de \SI{0,5}{m}, calibrado a \SI{20}{\celsius}, funciona un día a
\SI{40}{\celsius}. ¿Se adelanta o se atrasa? ¿Cuánto? ($\alpha_l=\SI{17e-6}{/K}$.)
\end{problema}

\begin{resolucion}
\idea El periodo $\mathcal T=2\pi\sqrt{l/g}$ crece con la longitud. El reloj cuenta oscilaciones: si cada
una dura más, marca menos tiempo del real y se \emph{atrasa}. Método M3.

\paso{Paso 1: longitud a \SI{40}{\celsius}} Con $\Delta T=\Delta t=\SI{20}{K}$ y \eqref{eq:DL},
$l=l_0(1+\alpha_l\Delta T)$, donde $\alpha_l\Delta T=\num{3,4e-4}\ll1$.

\paso{Paso 2: variación relativa del periodo} En el cociente se cancelan $2\pi$ y $g$; con la
aproximación binomial $(1+x)^{1/2}\simeq1+\tfrac12x$:
\[
  \frac{\mathcal T}{\mathcal T_0}=\sqrt{\frac{l}{l_0}}=\sqrt{1+\alpha_l\Delta T}\simeq1+\tfrac12\alpha_l\Delta T
  \quad\Longrightarrow\quad
  \frac{\Delta\mathcal T}{\mathcal T_0}\simeq\tfrac12\alpha_l\Delta T=\tfrac12\cdot\num{17e-6}\cdot20=\num{1,70e-4}.
\]

\paso{Paso 3: lectura del reloj} En un tiempo real $t$ el péndulo da $N=t/\mathcal T$ oscilaciones, y el
reloj, calibrado para $\mathcal T_0$, marca $N\mathcal T_0=t\,\mathcal T_0/\mathcal T$. El retraso es
\[
  t-t\,\frac{\mathcal T_0}{\mathcal T}=t\Bigl(1-\frac{\mathcal T_0}{\mathcal T}\Bigr)\simeq t\,\frac{\Delta\mathcal T}{\mathcal T_0},
\]
donde se ha usado $\mathcal T_0/\mathcal T=(1+\Delta\mathcal T/\mathcal T_0)^{-1}\simeq1-\Delta\mathcal T/\mathcal T_0$.

\paso{Paso 4: retraso diario} Con $t=\SI{1}{día}=\SI{86400}{s}$:
\[
  \SI{86400}{s}\cdot\num{1,70e-4}=\SI{14,7}{s}\quad\Longrightarrow\quad\boxed{\,\text{se atrasa}\ \simeq\SI{14,7}{s/día}\,}
\]
\end{resolucion}

\begin{observacion}[La longitud no interviene]
El resultado no depende de $l_0=\SI{0,5}{m}$ ni de $g$: solo importa la variación \emph{relativa},
$\tfrac12\alpha_l\Delta T$. Es el rasgo general de los coeficientes de respuesta. Como el retraso es
lineal en $\Delta T$, por debajo de \SI{20}{\celsius} el reloj se adelanta (figura~\ref{fig:p15}).
\end{observacion}

\begin{figure}[htbp]
\centering
\begin{tikzpicture}
\begin{axis}[graf,width=7.6cm,height=5cm,xmin=0,xmax=42,ymin=-17,ymax=17,
    xtick={0,10,20,30,40},ytick={-14.7,0,14.7},yticklabels={$-14{,}7$,0,$14{,}7$},
    xlabel={$t$ (\si{\celsius})},ylabel={retraso (s/día)},
    extra y ticks={0},extra y tick style={grid=major,grid style={regla}}]
  \addplot[navy,domain=0:42,samples=2] {0.7344*(x-20)};
  \fill[rojo] (axis cs:40,14.69) circle (1.8pt);
  \fill[acento] (axis cs:20,0) circle (1.8pt) node[below right,font=\scriptsize] {calibración};
  \node[font=\scriptsize,rojo,anchor=east] at (axis cs:30,12) {\SI{40}{\celsius}: se atrasa};
  \node[font=\scriptsize,acento,anchor=west] at (axis cs:2,-12) {se adelanta};
\end{axis}
\end{tikzpicture}
\caption{Problema 1.5. Retraso diario $\simeq\tfrac12\alpha_l(t-\SI{20}{\celsius})\cdot\SI{86400}{s}$
de un péndulo de cobre calibrado a \SI{20}{\celsius}.}
\label{fig:p15}
\end{figure}

% ---------------------------------------------------------------------
\subsection{Problema 1.6. Dilatación de un puente}
% ---------------------------------------------------------------------
\begin{problema}{1.6}{Puente de acero}
Un puente de acero continuo mide \SI{100}{m}. ¿Cuánto varía su longitud entre \SI{-10}{\celsius} y
\SI{50}{\celsius}? ($\alpha_l=\SI{11e-6}{/K}$.)
\end{problema}

\begin{resolucion}
\idea Método M3 con $\alpha_l\Delta T\ll1$.

\paso{Paso 1: salto térmico} Un intervalo Celsius es igual a uno Kelvin:
$\Delta T=\Delta t=50-(-10)=\SI{60}{K}$.

\paso{Paso 2: validez de la aproximación lineal}
$\alpha_l\Delta T=\num{11e-6}\cdot60=\num{6,6e-4}\ll1$, así que \eqref{eq:DL} es aplicable.

\paso{Paso 3: cálculo}
\[
  \Delta L\simeq\alpha_lL_0\,\Delta T=\SI{11e-6}{/K}\cdot\SI{100}{m}\cdot\SI{60}{K}=\SI{0,066}{m}
  \quad\Longrightarrow\quad\boxed{\,\Delta L\simeq\SI{6,6}{cm}\,}
\]

\paso{Paso 4: error de la aproximación} El resultado exacto \eqref{eq:Lexp} es
$\Delta L=L_0(\e^{\alpha_l\Delta T}-1)$; frente a \eqref{eq:DL} añade el término
$\tfrac12L_0(\alpha_l\Delta T)^2\simeq\SI{22}{\micro m}$, del todo despreciable frente a \SI{6,6}{cm}.
\end{resolucion}

% ---------------------------------------------------------------------
\subsection{Problema 1.7. Mercurio calentado a volumen constante}
% ---------------------------------------------------------------------
\begin{problema}{1.7}{Presión a volumen fijo}
Calcúlese el aumento de presión de una masa de mercurio calentada a volumen constante de
\SI{0}{\celsius} a \SI{1}{\celsius}, con $\alpha=\SI{18e-5}{/K}$ y $\kappa_T=\SI{39e-12}{/Pa}$ constantes.
\end{problema}

\begin{resolucion}
\idea A volumen fijo, $\dd\ln V=0$ en \eqref{eq:dlnV}.

\paso{Paso 1: condición de volumen fijo}
\[
  0=\dd\ln V=\alpha\,\dd T-\kappa_T\,\dd p\quad\Longrightarrow\quad\dd p=\frac{\alpha}{\kappa_T}\,\dd T.
\]

\paso{Paso 2: integración} Con $\alpha$ y $\kappa_T$ constantes salen de la integral:
\[
  \int_{p_0}^{p_0+\Delta p}\dd p=\frac{\alpha}{\kappa_T}\int_{T_0}^{T_0+\Delta T}\dd T
  \quad\Longrightarrow\quad\Delta p=\frac{\alpha}{\kappa_T}\,\Delta T.
\]

\paso{Paso 3: cálculo} Con $\Delta T=\SI{1}{K}$ y $\SI{1}{bar}=\SI{e5}{Pa}$:
\[
  \Delta p=\frac{\SI{18e-5}{/K}}{\SI{39e-12}{/Pa}}\cdot\SI{1}{K}=\SI{4,615e6}{Pa}
  \quad\Longrightarrow\quad\boxed{\,\Delta p\simeq\SI{46,2}{bar}\,}
\]

\paso{Paso 4: comprobación en dos etapas} Calentando a $p$ constante, $\Delta V_1/V=\alpha\Delta T=\num{1,8e-4}$;
para recuperar el volumen hay que comprimir con $\kappa_T\Delta p=\num{1,8e-4}$, es decir,
$\Delta p=\num{1,8e-4}/\SI{39e-12}{/Pa}=\SI{4,6e6}{Pa}$\chk\ (figura~\ref{fig:dosetapas}).
\end{resolucion}

\begin{observacion}[Orden de magnitud]
Un solo kelvin produce unas 46 atmósferas: los líquidos son casi incompresibles ($\kappa_T$ minúsculo),
por lo que impedir su dilatación exige presiones enormes. De aquí el riesgo de un termómetro
sobrecalentado. Nótese que $\alpha/\kappa_T=p\beta=\pdi{p}{T}{V}$.
\end{observacion}

% ---------------------------------------------------------------------
\subsection{Problema 1.8. Coeficientes de la ecuación de Clausius}
% ---------------------------------------------------------------------
\begin{problema}{1.8}{Ecuación de Clausius}
Para la ecuación de estado $p(\Vm-b)=RT$, calcúlense $\alpha$ y $\kappa_T$.
\end{problema}

\begin{resolucion}
\idea Método M4: aquí $\Vm$ se despeja sin dificultad.

\paso{Paso 0: despeje} $\displaystyle\Vm-b=\frac{RT}{p}\ \Longrightarrow\ \Vm=\frac{RT}{p}+b=\frac{RT+bp}{p}$.

\paso{Paso 1: derivadas parciales} $b$ es constante; $RT/p$ es lineal en $T$ y va como $p^{-1}$:
\[
  \pd{\Vm}{T}{p}=\frac{R}{p},\qquad\pd{\Vm}{p}{T}=RT\cdot(-p^{-2})=-\frac{RT}{p^2}.
\]

\paso{Paso 2: coeficientes} Dividiendo por $\Vm=(RT+bp)/p$, es decir, multiplicando por $p/(RT+bp)$:
\[
  \alpha=\frac{R}{p}\cdot\frac{p}{RT+bp}\ \Longrightarrow\ \boxed{\,\alpha=\frac{R}{RT+bp}\,}\qquad
  \kappa_T=\frac{RT}{p^2}\cdot\frac{p}{RT+bp}\ \Longrightarrow\ \boxed{\,\kappa_T=\frac{RT}{p\,(RT+bp)}\,}
\]

\paso{Paso 3: en variables \texorpdfstring{$(T,\Vm)$}{(T,V)}} Con $p=RT/(\Vm-b)$,
$RT+bp=RT\bigl(1+\frac{b}{\Vm-b}\bigr)=\frac{RT\Vm}{\Vm-b}$, y por tanto
\[
  \alpha=\frac{\Vm-b}{T\,\Vm},\qquad\kappa_T=\frac{(\Vm-b)^2}{RT\,\Vm}.
\]

\paso{Paso 4: comprobación} De $p=RT/(\Vm-b)$: $\pdi{p}{T}{\Vm}=R/(\Vm-b)=p/T$, luego $\beta=1/T$, y
\[
  p\beta\kappa_T=p\cdot\frac1T\cdot\frac{RT}{p(RT+bp)}=\frac{R}{RT+bp}=\alpha\chk
\]
Con $b=0$, $\alpha=1/T$ y $\kappa_T=1/p$ (gas ideal).
\end{resolucion}

\begin{figure}[htbp]
\centering
\begin{tikzpicture}
\begin{axis}[graf,leyenda fuera,width=7.4cm,height=5cm,xmin=0,xmax=1,ymin=0.4,ymax=1.08,
    xlabel={$x=bp/RT$},ylabel={cociente con el gas ideal}]
  \addplot[navy,domain=0:1,samples=60] {1/(1+x)};
  \addplot[rojo,densely dashed,domain=0:1,samples=2] {1};
  \legend{$\alpha T=\kappa_Tp=1/(1+x)$,gas ideal}
\end{axis}
\end{tikzpicture}
\caption{Problema 1.8. El covolumen reduce por igual $\alpha$ y $\kappa_T$ respecto del gas ideal: el gas
de Clausius es menos dilatable y menos compresible. Las correcciones son de orden $bp/RT$.}
\label{fig:p18}
\end{figure}

% ---------------------------------------------------------------------
\subsection{Problema 1.9. Ecuación de estado a partir de \texorpdfstring{$\alpha$}{alpha} y \texorpdfstring{$\kappa_T$}{kappa} (I)}
% ---------------------------------------------------------------------
\begin{problema}{1.9}{$\alpha=3/T$, $\kappa_T=2/p$}
Determínese la ecuación de estado de un sistema hidrostático con $\alpha=3/T$ y $\kappa_T=2/p$.
\end{problema}

\begin{resolucion}
\idea Método M5 con $\dd\ln\Vm$: los coeficientes no contienen $\Vm$ y las variables están separadas.

\paso{Paso 0: planteamiento} Incógnita $\Vm=\Vm(T,p)$; datos $\alpha=3/T$, $\kappa_T=2/p$; sin límite físico.

\paso{Paso 1: diferencial total} Con \eqref{eq:derivadas}:
\[
  \dd\Vm=\pd{\Vm}{T}{p}\dd T+\pd{\Vm}{p}{T}\dd p=\Vm\alpha\,\dd T-\Vm\kappa_T\,\dd p.
\]

\paso{Paso 2: forma} $\alpha$ y $\kappa_T$ no contienen $\Vm$: se divide entre $\Vm$,
$\dd\ln\Vm=\alpha\,\dd T-\kappa_T\,\dd p$.

\paso{Paso 3: sustitución}
\[
  \dd\ln\Vm=\underbrace{\frac{3}{T}}_{M}\,\dd T\underbrace{-\frac{2}{p}}_{N}\,\dd p.
\]

\paso{Paso 4: exactitud}
\[
  \pd{M}{p}{T}=\frac{\partial}{\partial p}\Bigl(\frac3T\Bigr)=0,\qquad
  \pd{N}{T}{p}=\frac{\partial}{\partial T}\Bigl(-\frac2p\Bigr)=0\quad\Longrightarrow\quad\text{exacta.}
\]

\paso{Paso 5: integración (caso A)} $M=M(T)$ y $N=N(p)$:
\[
  \ln\Vm=3\int\frac{\dd T}{T}-2\int\frac{\dd p}{p}=3\ln T-2\ln p+\ln K.
\]

\paso{Paso 6: quitar logaritmos} Con $3\ln T=\ln T^3$, $2\ln p=\ln p^2$ y $\ln x-\ln y=\ln(x/y)$:
\[
  \ln\Vm=\ln\frac{K\,T^3}{p^2}\quad\Longrightarrow\quad\Vm=\frac{K\,T^3}{p^2}\quad\Longrightarrow\quad
  \boxed{\,p^2\Vm=\cte\cdot T^3\,}
\]
Si se hubiera escrito $+C$ en vez de $+\ln K$, tomando exponenciales
$\Vm=\e^{\ln(T^3/p^2)+C}=\e^{C}\,T^3/p^2$, con $K=\e^C$.

\paso{Paso 7: comprobación} Con $\Vm=KT^3p^{-2}$:
\[
  \alpha=\frac1\Vm\cdot\frac{3KT^2}{p^2}=\frac3T\chk\qquad
  \kappa_T=-\frac1\Vm\cdot\Bigl(-\frac{2KT^3}{p^3}\Bigr)=\frac2p\chk
\]
\end{resolucion}

\begin{figure}[htbp]
\centering
\begin{tikzpicture}
\begin{axis}[graf,leyenda fuera,width=7.6cm,height=5.2cm,xmin=0,xmax=2,ymin=0,ymax=4,
    xlabel={$T$ (u.\,arb.)},ylabel={$\Vm$ (u.\,arb.)}]
  \addplot[acentoclaro,domain=0:2,samples=60] {x^3/0.5^2/4};
  \addplot[acentomedio,domain=0:2,samples=60] {x^3/1^2/4};
  \addplot[navy,domain=0:2,samples=60] {x^3/1.5^2/4};
  \legend{$p=0{,}5$,$p=1$,$p=1{,}5$}
\end{axis}
\end{tikzpicture}
\caption{Problema 1.9. Isobaras $\Vm=KT^3/p^2$ ($K=1/4$). Su pendiente relativa,
$\frac1\Vm\frac{\dd\Vm}{\dd T}=3/T$, es el coeficiente $\alpha$ del enunciado.}
\label{fig:p19}
\end{figure}

% ---------------------------------------------------------------------
\subsection{Problema 1.10. Ecuación de estado a partir de \texorpdfstring{$\alpha$}{alpha} y \texorpdfstring{$\kappa_T$}{kappa} (II)}
% ---------------------------------------------------------------------
\begin{problema}{1.10}{Gas con corrección lineal en $p$}
Determínese la ecuación de estado de un gas con $\displaystyle\alpha=\frac{R}{p\Vm}$ y
$\displaystyle\kappa_T=\frac{RT}{p^2\Vm}+\frac{B}{2\Vm}$.
\end{problema}

\begin{resolucion}
\idea Método M5. Multiplicando por $\Vm$ desaparece el volumen de los coeficientes y se integra
$\dd\Vm$ directamente; las variables no están separadas (caso B).

\paso{Paso 0: planteamiento} Incógnita $\Vm=\Vm(T,p)$; datos $\alpha$, $\kappa_T$; condición: para $B\to0$
los coeficientes son los del gas ideal.

\paso{Paso 1: diferencial total} $\dd\Vm=\Vm\alpha\,\dd T-\Vm\kappa_T\,\dd p$.

\paso{Paso 2: forma} $\alpha$ y $\kappa_T$ llevan el factor $1/\Vm$: no se divide, se trabaja con $\dd\Vm$,
y los productos $\Vm\alpha$, $\Vm\kappa_T$ ya no contienen $\Vm$. (Con $\dd\ln\Vm$ quedaría $\Vm$ en el
segundo miembro y no podría integrarse.)

\paso{Paso 3: sustitución}
\[
  \dd\Vm=\underbrace{\frac{R}{p}}_{M}\,\dd T\underbrace{-\Bigl(\frac{RT}{p^2}+\frac{B}{2}\Bigr)}_{N}\,\dd p.
\]

\paso{Paso 4: exactitud} Con la regla del cociente ($f=R$ constante, $g=p$, $g'=1$) y derivando
$N$ término a término:
\[
  \pd{M}{p}{T}=\frac{0\cdot p-R\cdot1}{p^2}=-\frac{R}{p^2},\qquad
  \pd{N}{T}{p}=-\frac{R}{p^2}-0=-\frac{R}{p^2}.
\]
Coinciden: exacta.

\paso{Paso 5: integración (caso B)} $M=R/p$ depende de $p$, que multiplica a $\dd T$: no están separadas.

\emph{5.1. Integrar en $T$ con $p$ fija:} $\displaystyle\Vm=\int\frac Rp\,\dd T=\frac{RT}{p}+g(p)$.

\emph{5.2. Derivar respecto de $p$ e igualar a $N$:}
\[
  \pd{\Vm}{p}{T}=-\frac{RT}{p^2}+g'(p)\stackrel{!}{=}-\frac{RT}{p^2}-\frac B2
  \quad\Longrightarrow\quad g'(p)=-\frac B2.
\]
Los términos en $T$ se cancelan, como debe ocurrir.

\emph{5.3. Integrar $g'$:} $g(p)=-\frac B2\,p+\Vm_0$, y por tanto
$\displaystyle\Vm=\frac{RT}{p}-\frac B2\,p+\Vm_0$.

\paso{Paso 6: constante} Para $B\to0$ los coeficientes son los del gas ideal y debe recuperarse
$p\Vm=RT$; como $p\Vm=RT+p\Vm_0$, eso exige $\Vm_0=0$. Multiplicando por $p$:
\[
  \boxed{\,p\Vm=RT-\tfrac12Bp^2\,}
\]

\paso{Paso 7: comprobación} Con $\Vm=RT/p-\tfrac B2p$:
\[
  \alpha=\frac1\Vm\cdot\frac Rp=\frac{R}{p\Vm}\chk\qquad
  \kappa_T=-\frac1\Vm\Bigl(-\frac{RT}{p^2}-\frac B2\Bigr)=\frac{RT}{p^2\Vm}+\frac{B}{2\Vm}\chk
\]
\end{resolucion}

\begin{figure}[htbp]
\centering
\begin{tikzpicture}
\begin{axis}[graf,leyenda fuera,width=7.6cm,height=5.2cm,xmin=0,xmax=3,ymin=0,ymax=2.3,
    xlabel={$p$ (u.\,arb.)},ylabel={$p\Vm$ (u.\,arb.)}]
  \addplot[acentoclaro,domain=0:3,samples=40] {1-0.1*x^2};
  \addplot[acentomedio,domain=0:3,samples=40] {1.5-0.1*x^2};
  \addplot[navy,domain=0:3,samples=40] {2-0.1*x^2};
  \addplot[rojo,densely dashed,domain=0:3,samples=2] {2};
  \addplot[rojo,densely dashed,domain=0:3,samples=2,forget plot] {1.5};
  \addplot[rojo,densely dashed,domain=0:3,samples=2,forget plot] {1};
  \legend{$RT=1$,$RT=1{,}5$,$RT=2$,gas ideal}
\end{axis}
\end{tikzpicture}
\caption{Problema 1.10. Diagrama de Amagat: $p\Vm=RT-\tfrac12Bp^2$ ($B=0{,}2$). Las isotermas son
parábolas que parten de las rectas horizontales del gas ideal en $p\to0$.}
\label{fig:p110}
\end{figure}

% ---------------------------------------------------------------------
\subsection{Problema 1.11. Ecuación de estado a partir de \texorpdfstring{$\beta$}{beta} y \texorpdfstring{$\kappa_T$}{kappa}}
% ---------------------------------------------------------------------
\begin{problema}{1.11}{$\beta=aT/(bp^2)$, $\kappa_T=bp/\Vm$}
Determínese la ecuación de estado de un sistema hidrostático con $\displaystyle\beta=\frac{aT}{bp^2}$ y
$\displaystyle\kappa_T=\frac{bp}{\Vm}$, siendo $a$ y $b$ constantes.
\end{problema}

\begin{resolucion}
\idea Método M5: completar el par con $\alpha=p\beta\kappa_T$ y usar $\dd\Vm$.

\paso{Paso 0: planteamiento} Incógnita $\Vm=\Vm(T,p)$. La diferencial de $\Vm$ necesita $\alpha$ y
$\kappa_T$; falta $\alpha$.

\paso{Paso 0 bis: tercer coeficiente} De la relación cíclica \eqref{eq:apbk}:
\[
  \alpha=p\,\beta\,\kappa_T=p\cdot\frac{aT}{bp^2}\cdot\frac{bp}{\Vm}=\frac{aT}{\Vm}.
\]

\paso{Paso 1--2: diferencial y forma} Los dos coeficientes llevan $1/\Vm$: se trabaja con $\dd\Vm$,
\[
  \dd\Vm=\Vm\alpha\,\dd T-\Vm\kappa_T\,\dd p.
\]

\paso{Paso 3: sustitución}
\[
  \dd\Vm=\underbrace{aT}_{M}\,\dd T\underbrace{-\,bp}_{N}\,\dd p.
\]

\paso{Paso 4: exactitud}
\[
  \pd{M}{p}{T}=\frac{\partial(aT)}{\partial p}=0,\qquad\pd{N}{T}{p}=\frac{\partial(-bp)}{\partial T}=0
  \quad\Longrightarrow\quad\text{exacta.}
\]

\paso{Paso 5: integración (caso A)} $M=M(T)$, $N=N(p)$:
\[
  \Vm=\int aT\,\dd T-\int bp\,\dd p\quad\Longrightarrow\quad
  \boxed{\,\Vm=\frac{aT^2}{2}-\frac{bp^2}{2}+\cte\,}
\]

\paso{Paso 6: constante} El enunciado no da ningún límite ni estado de referencia: queda libre.

\paso{Paso 7: comprobación}
\[
  \alpha=\frac{aT}{\Vm}\chk\qquad\kappa_T=-\frac{1}{\Vm}(-bp)=\frac{bp}{\Vm}\chk\qquad
  \beta=\frac{\alpha}{p\kappa_T}=\frac{aT/\Vm}{p\cdot bp/\Vm}=\frac{aT}{bp^2}\chk
\]
\end{resolucion}

\begin{figure}[htbp]
\centering
\begin{tikzpicture}
\begin{axis}[graf,leyenda fuera,width=7.2cm,height=5.4cm,xmin=0,xmax=3,ymin=0,ymax=3,
    xlabel={$T$ (u.\,arb.)},ylabel={$p$ (u.\,arb.)}]
  \addplot[acentoclaro,domain=1:3,samples=80] {sqrt(x^2-1)};
  \addplot[acentomedio,domain=1.4142:3,samples=80] {sqrt(x^2-2)};
  \addplot[navy,domain=2:3,samples=80] {sqrt(x^2-4)};
  \addplot[rojo,densely dashed,domain=0:3,samples=2] {x};
  \legend{$\Vm-\cte=0{,}5$,$\Vm-\cte=1$,$\Vm-\cte=2$,$\Vm-\cte=0$}
\end{axis}
\end{tikzpicture}
\caption{Problema 1.11. Isócoras en el plano $(T,p)$ con $a=b=1$: son las hipérbolas
$T^2-p^2=2(\Vm-\cte)$. Su pendiente, $\dd p/\dd T=aT/(bp)=p\beta$, es la que fija $\beta$.}
\label{fig:p111}
\end{figure}

% ---------------------------------------------------------------------
\subsection{Problema 1.12. Ecuación de estado a partir de \texorpdfstring{$\alpha$}{alpha} y \texorpdfstring{$\beta$}{beta}}
% ---------------------------------------------------------------------
\begin{problema}{1.12}{$\alpha$ y $\beta$ con denominador común}
Determínese la ecuación de estado de un fluido con $\displaystyle\alpha=\frac{a}{A+aT-bp}$ y
$\displaystyle\beta=\frac{a}{bp}$, siendo $A$, $a$ y $b$ constantes.
\end{problema}

\begin{resolucion}
\idea Método M5. Tras hallar $\kappa_T$, la diferencial es la logarítmica de $A+aT-bp$. Se resuelve
primero por el caso B, \emph{sin} ver el atajo, con la regla del cociente; el atajo se comenta al final.

\paso{Paso 0: planteamiento} Incógnita $\Vm=\Vm(T,p)$. Abreviamos $D\equiv A+aT-bp$, con
$\partial_TD=a$ y $\partial_pD=-b$.

\paso{Paso 0 bis: tercer coeficiente} De $\alpha=p\beta\kappa_T$, $\kappa_T=\alpha/(p\beta)$:
\[
  \kappa_T=\frac{a/D}{p\cdot\dfrac{a}{bp}}=\frac{a/D}{a/b}=\frac{a\cdot b}{D\cdot a}=\frac{b}{D}.
\]

\paso{Paso 1--2: diferencial y forma} $\alpha$ y $\kappa_T$ no contienen $\Vm$: se usa
$\dd\ln\Vm=\alpha\,\dd T-\kappa_T\,\dd p$.

\paso{Paso 3: sustitución}
\[
  \dd\ln\Vm=\underbrace{\frac{a}{D}}_{M}\,\dd T+\underbrace{\frac{-b}{D}}_{N}\,\dd p.
\]

\paso{Paso 4: exactitud (regla del cociente)} En ambos casos el numerador es constante ($f'=0$) y
$g=D$:
\[
  \pd{M}{p}{T}=\frac{0\cdot D-a\cdot(-b)}{D^2}=\frac{ab}{D^2},\qquad
  \pd{N}{T}{p}=\frac{0\cdot D-(-b)\cdot a}{D^2}=\frac{ab}{D^2}.
\]
Coinciden: exacta.

\paso{Paso 5: integración (caso B)} $M=a/D$ depende de $T$ y de $p$: no están separadas.

\emph{5.1. Integrar en $T$ con $p$ fija.} $A-bp\equiv c$ es constante; con $u=c+aT$, $\dd u=a\,\dd T$
(segunda primitiva de \eqref{eq:primitivas}):
\[
  \ln\Vm=\int\frac{a\,\dd T}{c+aT}+g(p)=\int\frac{\dd u}{u}+g(p)=\ln D+g(p)\qquad(D>0).
\]

\emph{5.2. Derivar respecto de $p$ e igualar a $N$.} $\partial_p\ln D=\partial_pD/D=-b/D$:
\[
  -\frac bD+g'(p)\stackrel{!}{=}-\frac bD\quad\Longrightarrow\quad g'(p)=0.
\]

\emph{5.3. Integrar $g'$.} $g$ es constante: $g=\ln K$.

\paso{Paso 6: quitar logaritmos}
\[
  \ln\Vm=\ln D+\ln K=\ln(KD)\quad\Longrightarrow\quad\boxed{\,\Vm=K\,(A+aT-bp)\,}
\]
Sin límite físico en el enunciado, $K$ queda libre.

\paso{Paso 7: comprobación} Con $\Vm=KD$, $\partial_T\Vm=Ka$ y $\partial_p\Vm=-Kb$:
\[
  \alpha=\frac{Ka}{KD}=\frac aD\chk\qquad\kappa_T=-\frac{-Kb}{KD}=\frac bD,\qquad
  \beta=\frac{\alpha}{p\kappa_T}=\frac{a/D}{pb/D}=\frac{a}{bp}\chk
\]
\end{resolucion}

\begin{observacion}[El atajo]
Que salga $g'=0$ indica que toda la dependencia estaba ya en $\ln D$. En efecto, los coeficientes
comparten el denominador $D$ y sus numeradores son $\partial_TD=a$ y $\partial_pD=-b$:
\[
  \dd\ln\Vm=\frac{a\,\dd T-b\,\dd p}{A+aT-bp}=\frac{\dd D}{D}=\dd\ln(A+aT-bp),
\]
que se integra en una línea. Las isócoras son rectas paralelas en el plano $(T,p)$ (figura~\ref{fig:p112}).
\end{observacion}

\begin{figure}[htbp]
\centering
\begin{tikzpicture}
\begin{axis}[graf,leyenda fuera,width=7.2cm,height=5.2cm,xmin=0,xmax=3,ymin=0,ymax=3.2,
    xlabel={$T$ (u.\,arb.)},ylabel={$p$ (u.\,arb.)}]
  \addplot[acentoclaro,domain=0:2.2,samples=2] {1+x-0.5};
  \addplot[acentomedio,domain=0:2.2,samples=2] {1+x-1};
  \addplot[navy,domain=0.5:3,samples=2] {1+x-1.5};
  \legend{$\Vm/K=0{,}5$,$\Vm/K=1$,$\Vm/K=1{,}5$}
\end{axis}
\end{tikzpicture}
\caption{Problema 1.12. Isócoras $p=(A+aT-\Vm/K)/b$ con $A=a=b=1$: rectas paralelas de pendiente
$a/b$. Al aumentar $\Vm$ se desplazan hacia abajo.}
\label{fig:p112}
\end{figure}

% ---------------------------------------------------------------------
\subsection{Problema 1.13. ¿Se recupera el gas ideal?}
% ---------------------------------------------------------------------
\begin{problema}{1.13}{Coeficientes con $\Vm-A$}
Un sistema hidrostático tiene $\displaystyle\alpha=\frac{\Vm-A}{2\Vm T}$ y
$\displaystyle\kappa_T=\frac{\Vm-A}{p\Vm}$ ($A$ constante). Determínese la ecuación térmica de estado.
¿Puede recuperarse el comportamiento de gas ideal?
\end{problema}

\begin{resolucion}
\idea Método M5. Al multiplicar por $\Vm$ aparece $\Vm-A$ en ambos términos: se integra
$\dd\ln(\Vm-A)$.

\paso{Paso 0: planteamiento} Incógnita $\Vm(T,p)$; datos $\alpha$, $\kappa_T$; se pregunta por el límite
$\Vm\gg A$.

\paso{Paso 1: diferencial total}
\[
  \dd\Vm=\Vm\alpha\,\dd T-\Vm\kappa_T\,\dd p=\frac{\Vm-A}{2T}\,\dd T-\frac{\Vm-A}{p}\,\dd p.
\]

\paso{Paso 2: forma} Los coeficientes contienen $\Vm-A$: se divide entre $\Vm-A$. Como $A$ es constante,
$\dd(\Vm-A)=\dd\Vm$ y el primer miembro es una diferencial logarítmica:
\[
  \frac{\dd\Vm}{\Vm-A}=\dd\ln(\Vm-A)=\underbrace{\frac{1}{2T}}_{M}\,\dd T\underbrace{-\frac1p}_{N}\,\dd p.
\]

\paso{Paso 3: exactitud} $\partial_p(1/2T)=0=\partial_T(-1/p)$: exacta.

\paso{Paso 4: integración (caso A)}
\[
  \ln(\Vm-A)=\frac12\ln T-\ln p+\ln K=\ln\frac{K\sqrt T}{p}
  \quad\Longrightarrow\quad\boxed{\,p\,(\Vm-A)=K\sqrt T\,}
\]

\paso{Paso 5: comprobación} De $\Vm=A+K\sqrt T/p$:
\[
  \partial_T\Vm=\frac{K}{2p\sqrt T}=\frac{\Vm-A}{2T}\ \Rightarrow\ \alpha=\frac{\Vm-A}{2\Vm T}\chk\qquad
  \partial_p\Vm=-\frac{K\sqrt T}{p^2}=-\frac{\Vm-A}{p}\ \Rightarrow\ \kappa_T=\frac{\Vm-A}{p\Vm}\chk
\]

\paso{Paso 6: límite} Para $\Vm\gg A$, $p\Vm\to K\sqrt T$, $\alpha\to1/(2T)$ y $\kappa_T\to1/p$. El gas ideal
exige $p\Vm=RT$ y $\alpha=1/T$: ninguna elección de $K$ convierte $K\sqrt T$ en $RT$ para todo $T$.
\textbf{No} se recupera el gas ideal en ningún límite, porque la dependencia en $T$ es $T^{1/2}$ y no $T$.
\end{resolucion}

\begin{observacion}[Lectura del enunciado]
Si el coeficiente fuera $\alpha=(\Vm-A)/(\Vm T)$, el mismo procedimiento daría
$\ln(\Vm-A)=\ln T-\ln p+\ln K$, es decir, $p(\Vm-A)=KT$: la ecuación de Clausius, que sí tiende al gas
ideal para $\Vm\gg A$ con $K=R$ (figura~\ref{fig:p113}). Conviene comprobar siempre el límite:
$\alpha\to1/T$ y $\kappa_T\to1/p$.
\end{observacion}

\begin{figure}[htbp]
\centering
\begin{tikzpicture}
\begin{axis}[graf,leyenda fuera,width=7.4cm,height=5cm,xmin=1,xmax=20,ymin=0,ymax=1.15,
    xlabel={$\Vm/A$},ylabel={$\alpha T$}]
  \addplot[navy,domain=1:20,samples=80] {(1-1/x)/2};
  \addplot[acento,domain=1:20,samples=80] {1-1/x};
  \addplot[rojo,densely dashed,domain=1:20,samples=2] {1};
  \legend{problema 1.13,Clausius,gas ideal}
\end{axis}
\end{tikzpicture}
\caption{Problema 1.13. $\alpha T$ frente a $\Vm/A$: con el enunciado dado tiende a $1/2$, no a $1$; con
$\alpha=(\Vm-A)/(\Vm T)$ (Clausius) sí tiende al valor ideal.}
\label{fig:p113}
\end{figure}

% ---------------------------------------------------------------------
\subsection{Problema 1.14. Ecuación de Dieterici}
% ---------------------------------------------------------------------
\begin{problema}{1.14}{$p(T,\Vm)$ a partir de $\kappa_T$ y $\beta$}
Un gas tiene $\displaystyle\kappa_T=\frac{RT\Vm(\Vm-b)}{p\,[RT\Vm^2-a(\Vm-b)]}$ y
$\displaystyle\beta=\frac{RT\Vm+a}{RT^2\Vm}$. Hállese $p(T,\Vm)$ sabiendo que para $\Vm$ grande el gas se
comporta como ideal.
\end{problema}

\begin{resolucion}
\idea Método M5 para la incógnita $p(T,\Vm)$: se diferencia $p$, no $\Vm$, y se usa
$\dd\ln p=\beta\,\dd T-\dd\Vm/(p\kappa_T\Vm)$, porque $\beta$ y $p\kappa_T$ no contienen $p$.

\paso{Paso 0: planteamiento} Incógnita $p=p(T,\Vm)$; datos $\beta$, $\kappa_T$; condición $p\Vm\to RT$
para $\Vm\to\infty$.

\paso{Paso 1: diferencial total} La primera derivada es $p\beta$ por definición; la segunda está
invertida respecto de la de $\kappa_T$ y se obtiene con la recíproca (ambas a $T$ constante):
\[
  \dd p=\pd{p}{T}{\Vm}\dd T+\pd{p}{\Vm}{T}\dd\Vm,\qquad
  \pd{p}{\Vm}{T}=\frac{1}{\pd{\Vm}{p}{T}}=-\frac{1}{\Vm\kappa_T}
  \quad\Longrightarrow\quad\dd p=p\beta\,\dd T-\frac{\dd\Vm}{\Vm\kappa_T}.
\]

\paso{Paso 2: forma} $\beta$ no contiene $p$ y $\kappa_T$ lleva el factor $1/p$, de modo que $p\kappa_T$
tampoco lo contiene. Dividiendo entre $p$:
\[
  \dd\ln p=\beta\,\dd T-\frac{\dd\Vm}{p\kappa_T\Vm}.
\]

\paso{Paso 3: sustitución y simplificación} Se separa cada coeficiente en fracciones simples:
\[
  \beta=\frac{RT\Vm}{RT^2\Vm}+\frac{a}{RT^2\Vm}=\frac1T+\frac{a}{RT^2\Vm};\qquad
  p\kappa_T\Vm=\frac{RT\Vm^2(\Vm-b)}{RT\Vm^2-a(\Vm-b)},
\]
\[
  \frac{1}{p\kappa_T\Vm}=\frac{RT\Vm^2-a(\Vm-b)}{RT\Vm^2(\Vm-b)}
  =\frac{RT\Vm^2}{RT\Vm^2(\Vm-b)}-\frac{a(\Vm-b)}{RT\Vm^2(\Vm-b)}=\frac{1}{\Vm-b}-\frac{a}{RT\Vm^2}.
\]
Por tanto
\[
  \dd\ln p=\underbrace{\Bigl(\frac1T+\frac{a}{RT^2\Vm}\Bigr)}_{M}\dd T
  +\underbrace{\Bigl(-\frac{1}{\Vm-b}+\frac{a}{RT\Vm^2}\Bigr)}_{N}\dd\Vm.
\]

\paso{Paso 4: exactitud (regla del cociente)} Los términos $1/T$ y $-1/(\Vm-b)$ no dependen de la otra
variable y dan derivada nula. Para los otros dos, $f=a$ ($f'=0$):
\[
  \pd{M}{\Vm}{T}=\frac{0-a\cdot RT^2}{(RT^2\Vm)^2}=-\frac{a}{RT^2\Vm^2},\qquad
  \pd{N}{T}{\Vm}=\frac{0-a\cdot R\Vm^2}{(RT\Vm^2)^2}=-\frac{a}{RT^2\Vm^2}.
\]
Coinciden: exacta.

\paso{Paso 5: integración (caso B)} $M$ depende de $T$ y de $\Vm$; la función auxiliar es ahora $g(\Vm)$.

\emph{5.1. Integrar en $T$ con $\Vm$ fijo.} Con $\int T^{-2}\dd T=-T^{-1}$:
\[
  \ln p=\int\Bigl(\frac1T+\frac{a}{R\Vm}\,T^{-2}\Bigr)\dd T+g(\Vm)=\ln T-\frac{a}{RT\Vm}+g(\Vm).
\]

\emph{5.2. Derivar respecto de $\Vm$ e igualar a $N$.} Por la regla del cociente ($f=-a$, $g=RT\Vm$,
$g'=RT$),
\[
  \frac{\partial}{\partial\Vm}\Bigl(\frac{-a}{RT\Vm}\Bigr)=\frac{0\cdot RT\Vm-(-a)\cdot RT}{(RT\Vm)^2}=\frac{a}{RT\Vm^2},
\]
de modo que
\[
  \frac{a}{RT\Vm^2}+g'(\Vm)\stackrel{!}{=}-\frac{1}{\Vm-b}+\frac{a}{RT\Vm^2}
  \quad\Longrightarrow\quad g'(\Vm)=-\frac{1}{\Vm-b}.
\]

\emph{5.3. Integrar $g'$.} Con $u=\Vm-b$: $g(\Vm)=-\ln(\Vm-b)+\ln K$. Así,
\[
  \ln p=\ln T-\ln(\Vm-b)+\ln K-\frac{a}{RT\Vm}=\ln\frac{KT}{\Vm-b}-\frac{a}{RT\Vm}.
\]

\paso{Paso 6: quitar logaritmos y fijar la constante} Queda un término que no es logaritmo; con
$\e^{\ln X+y}=X\e^y$:
\[
  p=\frac{KT}{\Vm-b}\,\exp\Bigl(-\frac{a}{RT\Vm}\Bigr).
\]
Para $\Vm\to\infty$: $\Vm-b\simeq\Vm$ y $\exp(-a/RT\Vm)\to\e^0=1$, luego $p\Vm\to KT$. El gas ideal exige
$K=R$:
\[
  \boxed{\,p=\frac{RT}{\Vm-b}\,\exp\Bigl(-\frac{a}{RT\Vm}\Bigr)\,}\qquad\text{(ecuación de Dieterici)}
\]

\paso{Paso 7: comprobación} De $\ln p=\ln R+\ln T-\ln(\Vm-b)-a/(RT\Vm)$:
\[
  \pd{\ln p}{T}{\Vm}=\frac1T+\frac{a}{RT^2\Vm}=\beta\chk\qquad
  \pd{\ln p}{\Vm}{T}=-\frac{1}{\Vm-b}+\frac{a}{RT\Vm^2}=N=-\frac{1}{p\kappa_T\Vm}\chk
\]
\end{resolucion}

\begin{observacion}[Compresibilidad negativa y punto crítico]\label{obs:kneg}
Como $\Vm>b$ y $p>0$, el signo de $\kappa_T$ es el de su denominador:
\[
  \kappa_T<0\iff RT\Vm^2-a(\Vm-b)<0\iff\pd{p}{\Vm}{T}>0.
\]
Un estado con $\kappa_T<0$ es \emph{mecánicamente inestable}: si el pistón de la figura~\ref{fig:piston}
se desplaza hacia fuera ($\delta\Vm>0$), la presión del gas \emph{sube} y lo empuja aún más; con
$\kappa_T>0$ baja y la perturbación se corrige. Equivalentemente, el equilibrio a $(T,\Vm)$ fijos es un
mínimo de la energía libre de Helmholtz, que exige
$\pdi{^2F}{\Vm^2}{T}=-\pdi{p}{\Vm}{T}=1/(\Vm\kappa_T)>0$.

\paso{Espinodales} Los extremos de la isoterma son las raíces de $RT\Vm^2-a\Vm+ab=0$:
\[
  \Vm_\pm=\frac{a\pm\sqrt{a^2-4abRT}}{2RT},
\]
reales si y solo si $a\geq4bRT$, es decir, $T\leq T_c$ con
\[
  T_c=\frac{a}{4Rb},\qquad\Vm_c=\frac{a}{2RT_c}=2b,\qquad p_c=\frac{RT_c}{\Vm_c-b}\,\e^{-a/(RT_c\Vm_c)}=\frac{a}{4\e^2b^2}.
\]
En $T_c$ las dos raíces se juntan en un punto de inflexión horizontal: es el \emph{punto crítico}. Por
debajo de $T_c$, entre $\Vm_-$ y $\Vm_+$ el fluido homogéneo es inestable y se separa en líquido y
vapor a la presión de vapor, que se obtiene con la \emph{construcción de Maxwell} (áreas iguales). Con
los valores de la figura~\ref{fig:p114}, $T_c\simeq\SI{304}{K}$, $\Vm_c\simeq\SI{0,093}{L/mol}$ y
$p_c\simeq\SI{73,9}{bar}$; a \SI{270}{K}, $\Vm_-\simeq\SI{0,069}{L/mol}$, $\Vm_+\simeq\SI{0,139}{L/mol}$ y la
presión de vapor es $\simeq\SI{52,2}{bar}$.
\end{observacion}

\begin{figure}[htbp]
\centering
\begin{tikzpicture}
\begin{axis}[graf,leyenda fuera,width=8.6cm,height=6.2cm,xmin=0.04,xmax=0.6,ymin=0,ymax=130,
    xtick={0.1,0.2,0.3,0.4,0.5,0.6},ytick={0,25,50,75,100,125},
    xlabel={$\Vm$ (L/mol)},ylabel={$p$ (bar)},restrict y to domain=-10:600,
    xticklabel style={/pgf/number format/fixed,/pgf/number format/precision=1}]
  \fill[rojo!9] (axis cs:0.0694,0) rectangle (axis cs:0.1391,130);
  \addplot[acentoclaro,domain=0.0485:0.6,samples=260] {0.08314*270/(x-0.0463)*exp(-4.68/(0.08314*270*x))};
  \addplot[acento,domain=0.0485:0.6,samples=260] {0.08314*303.94/(x-0.0463)*exp(-4.68/(0.08314*303.94*x))};
  \addplot[navy,domain=0.0485:0.6,samples=260] {0.08314*345/(x-0.0463)*exp(-4.68/(0.08314*345*x))};
  \addplot[rojo,densely dotted,domain=0.21:0.6,samples=80] {0.08314*345/x};
  \addplot[verde,densely dashed,thick,samples=2] coordinates {(0.0584,52.22) (0.1887,52.22)};
  \fill[verde] (axis cs:0.0584,52.22) circle (1.5pt);
  \fill[verde] (axis cs:0.1887,52.22) circle (1.5pt);
  \fill[rojo] (axis cs:0.0694,48.18) circle (1.5pt);
  \fill[rojo] (axis cs:0.1391,54.04) circle (1.5pt);
  \fill[acento] (axis cs:0.0926,73.86) circle (1.8pt) node[above right,font=\scriptsize] {C};
  \node[rojo,font=\scriptsize] at (axis cs:0.1043,20) {$\kappa_T<0$};
  \legend{$T=\SI{270}{K}<T_c$,$T=T_c\simeq\SI{304}{K}$,$T=\SI{345}{K}$,ideal \SI{345}{K},Maxwell \SI{270}{K}}
\end{axis}
\end{tikzpicture}
\caption{Problema 1.14. Isotermas de Dieterici con valores ilustrativos ($a=\SI{4,68}{bar.L^2/mol^2}$,
$b=\SI{0,0463}{L/mol}$, $T_c=a/4Rb$). A volúmenes grandes tienden a la isoterma ideal (punteada). Por
debajo de $T_c$ aparece el tramo sombreado con $\pdi{p}{\Vm}{T}>0$, es decir $\kappa_T<0$, entre las
espinodales (puntos rojos); la recta de Maxwell (verde) es la isoterma real de coexistencia
líquido--vapor. C es el punto crítico.}
\label{fig:p114}
\end{figure}

% ---------------------------------------------------------------------
\subsection{Problema 1.15. Coeficiente de dilatación de un gas del virial}
% ---------------------------------------------------------------------
\begin{problema}{1.15}{$\alpha$ en un estado dado}
Un gas tiene $\displaystyle\kappa_T=\frac{1}{RT\,(1/\Vm+2b/\Vm^2)}$ y $\displaystyle\beta=\frac1T$, con
$b=\SI{4e-2}{L/mol}$. Calcúlese $\alpha$ cuando un mol ocupa \SI{1}{L} a \SI{300}{K}, sabiendo que para
$\Vm$ grande el gas es ideal.
\end{problema}

\begin{resolucion}
\idea $\alpha=p\beta\kappa_T$ exige conocer $p$ en ese estado, luego antes hay que reconstruir
$p(T,\Vm)$ (método M5). Se integra primero la ecuación que no contiene la incógnita.

\paso{Paso 0: planteamiento} Incógnita $p(T,\Vm)$; datos $\beta$, $\kappa_T$; condición $p\Vm\to RT$ para
$\Vm\to\infty$; después, evaluar $\alpha$ en $\Vm=\SI{1}{L/mol}$, $T=\SI{300}{K}$.

\paso{Paso 1: diferencial total} Por \eqref{eq:dp}, con
$\frac{1}{\Vm\kappa_T}=\frac{RT}{\Vm}\bigl(\frac1\Vm+\frac{2b}{\Vm^2}\bigr)$:
\[
  \dd p=p\beta\,\dd T-\frac{\dd\Vm}{\kappa_T\Vm}
  =\underbrace{\frac pT}_{M}\,\dd T\underbrace{-RT\Bigl(\frac{1}{\Vm^2}+\frac{2b}{\Vm^3}\Bigr)}_{N}\,\dd\Vm.
\]

\paso{Paso 2: forma y orden de integración} Aquí $\kappa_T$ no contiene $p$, pero $p\beta=p/T$ sí:
$M$ depende de la incógnita. Por eso se integra primero $N$, que no la contiene, y después se usa
$M$ para fijar la función auxiliar (que ahora depende de $T$).

\paso{Paso 3: integración en \texorpdfstring{$\Vm$}{V} con \texorpdfstring{$T$}{T} fija} Con
$\int\Vm^{-2}\dd\Vm=-\Vm^{-1}$ y $\int\Vm^{-3}\dd\Vm=-\tfrac12\Vm^{-2}$:
\[
  p=-RT\Bigl(-\frac1\Vm-\frac{2b}{2\Vm^2}\Bigr)+f(T)=RT\Bigl(\frac1\Vm+\frac{b}{\Vm^2}\Bigr)+f(T).
\]

\paso{Paso 4: imponer la ecuación en \texorpdfstring{$T$}{T}} Derivando a $\Vm$ fijo e igualando a $M=p/T$:
\[
  R\Bigl(\frac1\Vm+\frac b{\Vm^2}\Bigr)+f'(T)\stackrel{!}{=}\frac pT=R\Bigl(\frac1\Vm+\frac b{\Vm^2}\Bigr)+\frac{f(T)}{T}
  \quad\Longrightarrow\quad f'(T)=\frac{f(T)}{T}.
\]
Separando variables, $\dd f/f=\dd T/T\Rightarrow\ln f=\ln T+\ln c$, es decir, $f=cT$.

\paso{Paso 5: constante} $p\Vm=RT(1+b/\Vm)+cT\Vm$. Para $\Vm\to\infty$ el primer término tiende a $RT$,
pero el segundo diverge salvo que $c=0$:
\[
  \boxed{\,p=\frac{RT}{\Vm}\Bigl(1+\frac{b}{\Vm}\Bigr)\,}
\]

\paso{Paso 6: comprobación}
\begin{align*}
  \pd{p}{T}{\Vm}&=\frac R\Vm\Bigl(1+\frac b\Vm\Bigr)=\frac pT &&\Longrightarrow\quad \beta=\frac1T\chk\\
  \pd{p}{\Vm}{T}&=-RT\Bigl(\frac1{\Vm^2}+\frac{2b}{\Vm^3}\Bigr) &&\Longrightarrow\quad
  \kappa_T=\frac{-1}{\Vm\,\pdi{p}{\Vm}{T}}=\frac{1}{RT(1/\Vm+2b/\Vm^2)}\chk
\end{align*}

\paso{Paso 7: coeficiente de dilatación} Multiplicando numerador y denominador por $\Vm^2$:
\[
  \alpha=p\beta\kappa_T=RT\Bigl(\frac1\Vm+\frac b{\Vm^2}\Bigr)\cdot\frac1T\cdot\frac{1}{RT(1/\Vm+2b/\Vm^2)}
  =\frac{\Vm+b}{T\,(\Vm+2b)}.
\]
Con $\Vm=\SI{1}{L/mol}$ y $b=\SI{0,04}{L/mol}$ (mismas unidades, el cociente es adimensional):
\[
  \alpha=\frac{1{,}04}{300\cdot1{,}08}\ \si{1/K}\quad\Longrightarrow\quad\boxed{\,\alpha\simeq\SI{3,21e-3}{/K}\,}
\]
\end{resolucion}

\begin{observacion}[Comparación con el gas ideal]
El gas ideal daría $\alpha=1/T=\SI{3,33e-3}{/K}$; la corrección relativa es
$\frac{\Vm+b}{\Vm+2b}-1=-\frac{b}{\Vm+2b}\simeq\SI{-3,7}{\percent}$, del orden de $b/\Vm$, como corresponde
a un desarrollo del virial truncado (figura~\ref{fig:p115}).
\end{observacion}

\begin{figure}[htbp]
\centering
\begin{tikzpicture}
\begin{axis}[graf,leyenda fuera,width=7.4cm,height=5cm,xmin=0.1,xmax=2,ymin=2.4,ymax=3.5,
    xlabel={$\Vm$ (L/mol)},ylabel={$\alpha$ ($10^{-3}$/K)}]
  \addplot[navy,domain=0.1:2,samples=80] {1000*(x+0.04)/(300*(x+0.08))};
  \addplot[rojo,densely dashed,domain=0.1:2,samples=2] {1000/300};
  \fill[acento] (axis cs:1,3.2099) circle (1.8pt) node[below,font=\scriptsize] {$3{,}21$};
  \legend{virial,gas ideal}
\end{axis}
\end{tikzpicture}
\caption{Problema 1.15. Coeficiente de dilatación del gas del virial a \SI{300}{K} frente al volumen
molar. El punto marca el estado del enunciado; para $\Vm\gg b$ se recupera $1/T$.}
\label{fig:p115}
\end{figure}

% ---------------------------------------------------------------------
\subsection{Problema 1.16. Cobre calentado y comprimido}
% ---------------------------------------------------------------------
\begin{problema}{1.16}{Volumen final con coeficientes medios}
Un kilogramo de cobre se calienta de \SI{300}{K} a \SI{500}{K} mientras su presión aumenta
\SI{1000}{bar}. Con $\alpha=\SI{51,7e-6}{/K}$, $\kappa_T=\SI{0,8065e-11}{/Pa}$ (valores medios) y
$v_0=\SI{111,5e-6}{m^3/kg}$, calcúlese el volumen final.
\end{problema}

\begin{resolucion}
\idea Método M3: integrar $\dd\ln V=\alpha\,\dd T-\kappa_T\,\dd p$ con coeficientes constantes. Como
$\ln V$ es función de estado, no importa en qué orden se calienta y se comprime (figura~\ref{fig:p116}).

\paso{Paso 0: datos en SI} $\Delta T=\SI{200}{K}$; $\Delta p=\SI{1000}{bar}=\SI{e8}{Pa}$;
$V_0=mv_0=\SI{111,5e-6}{m^3}$.

\paso{Paso 1: integración} Por el camino rectangular, \eqref{eq:const}:
\[
  \ln\frac{V}{V_0}=\int_{T_0}^{T_0+\Delta T}\alpha\,\dd T-\int_{p_0}^{p_0+\Delta p}\kappa_T\,\dd p=\alpha\,\Delta T-\kappa_T\,\Delta p.
\]

\paso{Paso 2: cálculo de cada término}
\[
  \alpha\Delta T=\SI{51,7e-6}{/K}\cdot\SI{200}{K}=\num{1,034e-2},\qquad
  \kappa_T\Delta p=\SI{0,8065e-11}{/Pa}\cdot\SI{e8}{Pa}=\num{8,065e-4},
\]
\[
  \ln\frac{V}{V_0}=\num{1,034e-2}-\num{8,065e-4}=\num{9,534e-3}.
\]

\paso{Paso 3: volumen final} Tomando exponenciales, $\e^{\num{9,534e-3}}=1{,}009580$:
\[
  V=V_0\,\e^{\num{9,534e-3}}=\SI{111,5e-6}{m^3}\cdot1{,}009580\quad\Longrightarrow\quad
  \boxed{\,V\simeq\SI{112,57e-6}{m^3}\,}
\]
\end{resolucion}

\begin{observacion}[Competencia de efectos]
La dilatación térmica ($+\SI{1,03}{\percent}$) domina sobre la compresión ($-\SI{0,08}{\percent}$). Como
$\ln(V/V_0)\ll1$, la aproximación lineal \eqref{eq:pequenos},
$\Delta V\simeq V_0(\alpha\Delta T-\kappa_T\Delta p)=\SI{1,063e-6}{m^3}$, da $V\simeq\SI{112,56e-6}{m^3}$: el
mismo resultado con cuatro cifras.
\end{observacion}

\begin{figure}[htbp]
\centering
\begin{tikzpicture}
\begin{axis}[graf,clip mode=individual,width=7.6cm,height=5.2cm,xmin=250,xmax=600,ymin=-200,ymax=1250,
    xtick={300,500},ytick={0,1000},xlabel={$T$ (K)},ylabel={$p-p_0$ (bar)}]
  \addplot[navy,thick,postaction={decorate},decoration={markings,mark=at position 0.55 with {\arrow{Stealth[length=5pt]}}}] coordinates {(300,0) (500,0)};
  \addplot[acento,thick,postaction={decorate},decoration={markings,mark=at position 0.55 with {\arrow{Stealth[length=5pt]}}}] coordinates {(500,0) (500,1000)};
  \addplot[gris,densely dotted] coordinates {(300,0) (300,1000) (500,1000)};
  \fill[rojo] (axis cs:300,0) circle (1.8pt);
  \fill[navy] (axis cs:500,0) circle (1.8pt);
  \fill[rojo] (axis cs:500,1000) circle (1.8pt);
  \node[font=\scriptsize,anchor=north east,fill=white,inner sep=1pt] at (axis cs:296,-30) {$111{,}50$};
  \node[font=\scriptsize,anchor=north west,fill=white,inner sep=1pt] at (axis cs:504,-30) {$112{,}66$};
  \node[font=\scriptsize,anchor=west,fill=white,inner sep=1pt] at (axis cs:508,1000) {$112{,}57$};
  \node[font=\scriptsize,navy,anchor=south] at (axis cs:400,30) {$+\alpha\Delta T$};
  \node[font=\scriptsize,acento,anchor=west] at (axis cs:508,500) {$-\kappa_T\Delta p$};
\end{axis}
\end{tikzpicture}
\caption{Problema 1.16. Camino rectangular en el plano $(T,p)$: calentamiento isobaro y compresión
isoterma. Junto a cada vértice, el volumen en $10^{-6}\,\mathrm{m^3}$. El camino punteado (comprimir
primero) da el mismo volumen final.}
\label{fig:p116}
\end{figure}

% ---------------------------------------------------------------------
\subsection{Parcial 2025/26, cuestión 1}
% ---------------------------------------------------------------------
\begin{problema}{P.1}{Parcial 2025/26}
Una sustancia hipotética tiene $\displaystyle\alpha=\frac{1}{T+A}$ y $\displaystyle\kappa_T=\frac1p$.
(a) Calcule la ecuación de estado $p=p(T,V)$. (b) Compruebe explícitamente la relación cíclica entre
$p$, $V$ y $T$. (c) En una expansión espontánea de este gas (presión inicial \SI{1}{bar}) hasta el
equilibrio con unos alrededores a $p_{\mathrm{ext}}=\SI{2}{bar}$, ¿pueden usarse estos coeficientes
para describir la curva del proceso?
\end{problema}

\begin{resolucion}
\idea (a) Método M5; (b) método M4 sobre el resultado; (c) concepto de proceso cuasiestático.

\paso{(a) Paso 0--2: planteamiento y forma} Aunque se pide $p(T,V)$, los datos son $\alpha$ y $\kappa_T$,
que no contienen $V$: se integra $\dd\ln V=\alpha\,\dd T-\kappa_T\,\dd p$ y al final se despeja $p$.

\paso{(a) Paso 3--4: sustitución y exactitud}
\[
  \dd\ln V=\underbrace{\frac{1}{T+A}}_{M}\,\dd T\underbrace{-\frac{1}{p}}_{N}\,\dd p,\qquad
  \pd{M}{p}{T}=0=\pd{N}{T}{p}\quad\Longrightarrow\quad\text{exacta.}
\]

\paso{(a) Paso 5--6: integración y despeje} Variables separadas (caso A); como $\dd(T+A)=\dd T$, la
primitiva es $\int\dd T/(T+A)=\ln(T+A)$:
\[
  \ln V=\ln(T+A)-\ln p+\ln K=\ln\frac{K(T+A)}{p}
  \quad\Longrightarrow\quad pV=K(T+A)\quad\Longrightarrow\quad\boxed{\,p=\frac{K\,(T+A)}{V}\,}
\]

\paso{(b) Relación cíclica} Despejando cada variable de $pV=K(T+A)$ y derivando:
\begin{alignat*}{3}
  p&=\frac{K(T+A)}{V} &\quad&\Longrightarrow\quad& \pd{p}{V}{T}&=-\frac{K(T+A)}{V^2}=-\frac pV,\\
  V&=\frac{K(T+A)}{p} &&\Longrightarrow& \pd{V}{T}{p}&=\frac Kp,\\
  T&=\frac{pV}{K}-A &&\Longrightarrow& \pd{T}{p}{V}&=\frac VK.
\end{alignat*}
Multiplicando las tres:
\[
  \pd{p}{V}{T}\pd{V}{T}{p}\pd{T}{p}{V}=\Bigl(-\frac pV\Bigr)\cdot\frac Kp\cdot\frac VK=-1\chk
\]
Equivalentemente, $\beta=\frac1p\pdi{p}{T}{V}=\frac1p\frac KV=\frac{1}{T+A}$ y
$p\beta\kappa_T=p\cdot\frac{1}{T+A}\cdot\frac1p=\frac{1}{T+A}=\alpha$\chk

\paso{(c) Proceso espontáneo} \textbf{No}. Una evolución espontánea con un desequilibrio finito de presión
($p_{\mathrm{ext}}-p_i=\SI{1}{bar}$) es irreversible y no cuasiestática: los estados intermedios no son de
equilibrio, $p$ y $T$ no están definidas uniformemente y el proceso no es una curva del espacio de
estados. Los coeficientes (y la ecuación de estado) solo relacionan los estados \emph{inicial y final},
ambos de equilibrio (figura~\ref{fig:pp1}).
\end{resolucion}

\begin{observacion}[Detalle del enunciado]
Con $p_{\mathrm{ext}}>p_i$ el proceso espontáneo sería en realidad una \emph{compresión}; la respuesta no
cambia, porque depende solo de que el proceso sea no cuasiestático. Tampoco la opción «solo si es
adiabática» la salva: ser adiabático no lo hace cuasiestático.
\end{observacion}

\begin{figure}[htbp]
\centering
\begin{tikzpicture}
\begin{axis}[graf,width=7.4cm,height=5.2cm,xmin=0.3,xmax=2.4,ymin=0.5,ymax=2.5,
    xtick=\empty,ytick={1,2},xlabel={$V$},ylabel={$p$ (bar)}]
  \addplot[acentoclaro,domain=0.45:2.4,samples=60] {1.6/x};
  \addplot[acentoclaro,domain=0.45:2.4,samples=60] {1.2/x};
  \fill[rojo] (axis cs:1.6,1) circle (2pt) node[below right,font=\scriptsize,black] {inicial};
  \fill[rojo] (axis cs:0.6,2) circle (2pt) node[above right,font=\scriptsize,black] {final};
  \draw[gris,densely dashed,-{Stealth[length=5pt]},decorate,decoration={snake,amplitude=2pt,segment length=9pt,post length=6pt}]
     (axis cs:1.55,1.05) -- (axis cs:0.66,1.93);
  \node[gris,font=\scriptsize,align=center] at (axis cs:1.55,1.9) {estados intermedios\\fuera de equilibrio};
\end{axis}
\end{tikzpicture}
\caption{Parcial, cuestión 1(c). Solo los estados inicial y final (puntos) son de equilibrio y están
sobre isotermas de la ecuación de estado; el proceso espontáneo (línea ondulada) no es una curva del
diagrama $p$--$V$ y los coeficientes no describen su trayectoria.}
\label{fig:pp1}
\end{figure}

% ---------------------------------------------------------------------
\subsection{Resumen de resultados}
% ---------------------------------------------------------------------
\begin{table}[H]
\centering
\caption{Resultados de los seminarios del Tema 1.}
\small
\renewcommand{\arraystretch}{1.3}
\begin{tabularx}{\textwidth}{lXc}
\toprule
Problema & Resultado & Método\\
\midrule
1.1 & $p_1(V_1-b)=p_2V_2^2/(V_2-b')$;\enspace $\theta_1=p_1(V_1-b)$, $\theta_2=p_2V_2^2/(V_2-b')$, $\theta_3=p_3V_3$ & M1\\
1.2 & $p^3V=KT^2$ & M5\\
1.3 & $p_2\simeq\SI{21,5}{kPa}$ & M2\\
1.4 & $\alpha\simeq1{,}18$ & M2\\
1.5 & se atrasa $\simeq\SI{14,7}{s/día}$ & M3\\
1.6 & $\Delta L\simeq\SI{6,6}{cm}$ & M3\\
1.7 & $\Delta p\simeq\SI{46,2}{bar}$ & M3\\
1.8 & $\alpha=R/(RT+bp)$,\enspace $\kappa_T=RT/[p(RT+bp)]$ & M4\\
1.9 & $p^2\Vm=\cte\cdot T^3$ & M5\\
1.10 & $p\Vm=RT-\tfrac12Bp^2$ & M5\\
1.11 & $\Vm=\tfrac12aT^2-\tfrac12bp^2+\cte$ & M5\\
1.12 & $\Vm=\cte\cdot(A+aT-bp)$ & M5\\
1.13 & $p(\Vm-A)=K\sqrt T$; no tiende al gas ideal & M5\\
1.14 & $p=\dfrac{RT}{\Vm-b}\,\e^{-a/(RT\Vm)}$ (Dieterici);\enspace $T_c=a/4Rb$ & M5\\
1.15 & $p=\frac{RT}{\Vm}\bigl(1+\frac b\Vm\bigr)$;\enspace $\alpha=\frac{\Vm+b}{T(\Vm+2b)}\simeq\SI{3,21e-3}{/K}$ & M5\\
1.16 & $V\simeq\SI{112,57e-6}{m^3}$ & M3\\
P.1 & $p=K(T+A)/V$; cíclica $=-1$; no (proceso no cuasiestático) & M5, M4\\
\bottomrule
\end{tabularx}
\end{table}

\begin{recuadro}{Lista de control antes de entregar}
\begin{itemize}[label={\color{acento}$\blacktriangleright$},leftmargin=1.4em,itemsep=4pt]
\item ¿Unidades homogéneas (\si{m^3}, \si{Pa}, \si{K})? En cocientes y en $\Delta T$ no hace falta convertir.
\item En M5, ¿se ha elegido la forma ($\dd X$ o $\dd\ln X$) en la que la incógnita no aparece en los coeficientes?
\item ¿Se ha comprobado la exactitud y, en el caso B, se cancelan los términos de la primera variable?
\item ¿Se ha fijado la constante con un límite físico, o se ha dicho que queda libre?
\item ¿El resultado recupera el gas ideal ($\alpha\to1/T$, $\kappa_T\to1/p$) cuando corresponde?
\item ¿Se cumple $\alpha=p\beta\kappa_T$ con los coeficientes obtenidos? ¿Y $\kappa_T>0$ (estabilidad)?
\item ¿El proceso es cuasiestático? Si no lo es, solo se pueden relacionar los estados extremos.
\end{itemize}
\end{recuadro}
